\documentclass[11pt,a4paper]{article}
\usepackage[T1]{fontenc}
\usepackage[utf8]{inputenc}
\usepackage[french]{babel}
\usepackage{geometry}
\usepackage{amsmath,amssymb}
\usepackage{booktabs}
\usepackage{hyperref}
\geometry{margin=2.3cm}
\hypersetup{colorlinks=true,linkcolor=black,urlcolor=blue}
\setlist{nosep}
\usepackage{enumitem}

\title{MarsWalk Navigator\\
\large Explication des fichiers, des champs et des pipelines}
\author{Dossier \texttt{nasa\_space\_2}}
\date{15 novembre 2026}

\begin{document}
\maketitle

\begin{abstract}
Le projet a deux fenêtres distinctes. La carte couleur sélectionnée est un contexte visuel.
Les trois chemins sont calculés sur une autre fenêtre de Gale, à partir d'une grille issue d'un MNT
(modèle numérique de terrain, DEM). Un chemin n'est jamais dessiné sur la carte couleur.
\end{abstract}

\section{Les deux zones}

\begin{center}
\begin{tabular}{@{}llp{6.2cm}@{}}
\toprule
Nom & Fichier & Rôle \\
\midrule
Contexte & \texttt{data/area.json} &
Carré d'environ \(2\,\mathrm{km}\times 2\,\mathrm{km}\).
Bornes \(137{,}37402^\circ\mathrm{E}\) à \(137{,}40788^\circ\mathrm{E}\),
\(-4{,}70125^\circ\) à \(-4{,}66751^\circ\). \\
Calcul de route & \texttt{data/route-area.json} &
Fenêtre boxwork séparée, environ \(900\,\mathrm{m}\times 677{,}5\,\mathrm{m}\).
Bornes \(137{,}29497^\circ\mathrm{E}\) à \(137{,}31023^\circ\mathrm{E}\),
\(-4{,}88177^\circ\) à \(-4{,}87032^\circ\). \\
\bottomrule
\end{tabular}
\end{center}

Ces deux rectangles ne se recouvrent pas. Une coordonnée lue sur la photo couleur
n'est pas un point de départ du calculateur.

\section{Pipeline A : carte de contexte}

Fichiers : \texttt{index.html}, \texttt{style.css}, \texttt{script.js},
\texttt{assets/maps/orbital.jpg}, \texttt{data/area.json}, \texttt{data/manifest.json}.

\begin{enumerate}
  \item Le navigateur charge \texttt{data/area.json} et \texttt{data/manifest.json}.
  \item L'image JPEG est affichée. Elle n'a pas de géoréférence interne.
  \item Chaque pixel est placé par interpolation linéaire des bornes.
  \item Un clic enregistre un départ ou une destination. Deux marqueurs apparaissent.
  \item La distance à vol d'oiseau est calculée. Aucune polyligne n'est dessinée.
  \item Le bouton \emph{Open DEM route calculation} ouvre \texttt{route-demo.html}.
\end{enumerate}

\subsection{Du clic à la longitude et à la latitude}

Soit \((x,y)\) la position du pointeur dans le rectangle de l'image, avec
\(x,y\in[0,1]\). L'origine \((0,0)\) est le coin nord-ouest.
\[
\lambda = \lambda_W + x\,(\lambda_E-\lambda_W),
\qquad
\varphi = \varphi_N - y\,(\varphi_N-\varphi_S).
\]
\(\lambda\) est la longitude est, \(\varphi\) la latitude planétocentrique.
Les quatre bornes viennent de \texttt{bounds} dans \texttt{area.json}.

\subsection{Distance à vol d'oiseau}

Ce n'est pas une route. C'est la longueur de l'arc entre les deux clics sur une sphère
de rayon moyen IAU \(R=3389{,}5\,\mathrm{km}\) :
\begin{align*}
a &= \sin^2\!\left(\frac{\Delta\varphi}{2}\right)
  + \cos\varphi_1\cos\varphi_2\sin^2\!\left(\frac{\Delta\lambda}{2}\right),\\
d &= 2R\arcsin\sqrt{a}.
\end{align*}
\(R\) n'est pas une mesure de ce site. La photo n'entre pas dans le calcul.

\section{Pipeline B : trois routes sur le MNT}

Fichiers : \texttt{route-demo.html}, \texttt{route-demo.css}, \texttt{route-demo.js},
\texttt{pathfinder.js}, \texttt{data/route-grid.js},
\texttt{assets/maps/route-terrain.jpg},
\texttt{data/dem/route\_dem\_source.tif}.

\begin{enumerate}
  \item Le MNT local a été réduit à une grille \(300\times 225\).
  \item La page charge cette grille, puis \texttt{pathfinder.js}.
  \item Les points d'exemple sont les cellules \((162,186)\) et \((216,150)\).
  \item Un clic sur l'image de terrain choisit la cellule la plus proche encore franchissable.
  \item Trois recherches A* à 8 voisins sont lancées, avec le même plafond de pente.
  \item Les polylignes sont dessinées seulement dans le SVG posé sur \texttt{route-terrain.jpg}.
\end{enumerate}

\subsection{Cellule et coordonnée}

Une cellule \((i,j)\) a son centre
\[
\lambda = \lambda_W + \frac{i+1/2}{W}(\lambda_E-\lambda_W),
\qquad
\varphi = \varphi_N - \frac{j+1/2}{H}(\varphi_N-\varphi_S),
\]
avec \(W=300\) et \(H=225\). Le pas de grille est \(c=3{,}005\,\mathrm{m}\).

\subsection{Voisinage et cases interdites}

Les huit voisins sont autorisés. Un pas diagonal n'est accepté que si les deux cases
adjacentes du carré sont elles aussi franchissables. Une case est bloquée si
\texttt{nodata} est vrai, ou si sa pente dépasse le plafond choisi \(s_{\max}\).
Le plafond par défaut est \(15^\circ\). Ce n'est pas une règle NASA : c'est un paramètre de l'application.

\subsection{Coût d'un pas}

Pour un pas de longueur \(\Delta s\) qui arrive sur une case de pente \(s\) et de score
de position \(I\in[0,1]\) :
\[
r=\frac{s}{s_{\max}},
\qquad
C=\Delta s\bigl(1+w_s r^2\bigr)+w_r\,\Delta s\,(1-I).
\]
Le pas orthogonal vaut \(c\), le pas diagonal vaut \(c\sqrt{2}\).
L'heuristique A* est la distance octile, donc admissible pour ce voisinage :
\[
h=c\bigl(\max(|\Delta i|,|\Delta j|)-\min(|\Delta i|,|\Delta j|)
+\sqrt{2}\,\min(|\Delta i|,|\Delta j|)\bigr).
\]

\subsection{Les trois jeux de poids}

\begin{center}
\begin{tabular}{@{}lcccl@{}}
\toprule
Option & \(w_s\) & \(w_r\) & Exemple à \(15^\circ\) & Sens \\
\midrule
Shortest & \(0\) & \(0\) & \(210{,}6\,\mathrm{m}\), pente max \(14{,}7^\circ\) &
distance seule, le plafond reste dur \\
Gentlest slope & \(12\) & \(0\) & \(221{,}2\,\mathrm{m}\), pente max \(7{,}5^\circ\) &
pénalité forte sur les cases raides \\
Ridge heuristic & \(1{,}5\) & \(8\) & \(226{,}4\,\mathrm{m}\), pente max \(8{,}0^\circ\) &
préfère un score topographique élevé \\
\bottomrule
\end{tabular}
\end{center}

Le score de crête est une heuristique de position dans le relief.
Ce n'est pas une valeur scientifique mesurée, ni une teneur minérale.

\subsection{Champs calculés sur un chemin}

Pour une suite de cellules, avec \(z\) l'altitude et \(s\) la pente :
\begin{align*}
D &= \sum \Delta s,\\
A &= \sum \max(z_{k}-z_{k-1},0),
\qquad
Descente = \sum \max(z_{k-1}-z_k,0),\\
s_{\mathrm{max}} &= \max s_k,
\qquad
\bar s = \frac{\sum s_k\Delta s}{D},
\qquad
\bar I = \frac{\sum I_k\Delta s}{D}.
\end{align*}
La fiche affiche aussi le nombre de cellules et le pas \(3{,}005\,\mathrm{m}\).
Elle dit que le résultat n'est pas une évaluation de sécurité.

\section{Champs des fichiers de données}

\subsection{\texttt{data/area.json}}

\begin{description}
\item[\texttt{coordinate\_system}] Longitude est, latitude planétocentrique, degrés décimaux.
\item[\texttt{center}] Centre enregistré du carré de contexte.
\item[\texttt{bounds}] \texttt{west}, \texttt{south}, \texttt{east}, \texttt{north}. Ce sont les seules bornes de la photo.
\item[\texttt{approximate\_size\_km}] Taille enregistrée, environ \(2{,}0\) par \(2{,}0\,\mathrm{km}\).
\item[\texttt{note}] Rappel : confirmer que la cible visible est bien dans le cadre avant une lecture scientifique.
\end{description}

\subsection{\texttt{data/route-area.json}}

Mêmes idées de bornes, plus :
\begin{description}
\item[\texttt{size\_m}] \(900{,}0\,\mathrm{m}\) est--ouest et \(677{,}5\,\mathrm{m}\) nord--sud.
\item[\texttt{grid}] \(300\times 225\), pas \(3{,}005\,\mathrm{m}\), mosaïque étiquetée \(1\,\mathrm{m}\).
\item[\texttt{retrieved\_on}] Copie du 15 novembre 2026. La date de récupération d'origine n'est pas connue.
\item[\texttt{note}] Fenêtre séparée. Les chemins sont des exemples de planification, pas des routes rover vérifiées.
\end{description}

\subsection{\texttt{data/route-grid.js}}

Objet \texttt{MARSWALK\_GRID}, généré, à ne pas éditer à la main.
\begin{description}
\item[\texttt{bounds}] Mêmes bornes que \texttt{route-area.json}.
\item[\texttt{width}, \texttt{height}] \(300\) et \(225\).
\item[\texttt{cell\_size\_m}] \(3{,}005\).
\item[\texttt{resolution\_m}] \(1\), étiquette de la mosaïque. La résolution réelle de chaque pixel de cette découpe n'est pas vérifiée.
\item[\texttt{coarse\_fill}] Indique si le remplissage grossier a été marqué.
\item[\texttt{width\_m}, \texttt{height\_m}] Emprise en mètres.
\item[\texttt{source}] USGS Astrogeology / projet MSL, mosaïque Gale, via NASA/JPL Mars Trek.
\item[\texttt{elev}] Altitude par case, en mètres.
\item[\texttt{slope}] Pente par case, en degrés, déjà dérivée du MNT.
\item[\texttt{interest}] Score de position topographique, entre \(0\) et \(1\).
\item[\texttt{nodata}] Case sans valeur : elle est infranchissable.
\end{description}

\subsection{Chaque entrée de \texttt{data/manifest.json}}

\begin{description}
\item[\texttt{id}] Identifiant stable : \texttt{orbital\_map}, \texttt{dem}, \texttt{route\_grid}, \texttt{route\_terrain}.
\item[\texttt{kind}] Nature : image, raster ou grille dérivée.
\item[\texttt{path}] Chemin relatif du fichier attendu.
\item[\texttt{source\_name}, \texttt{source\_url}] Nom et page de la source.
\item[\texttt{credit}] Citation à afficher.
\item[\texttt{license\_or\_use}] Contrainte d'usage, si elle est connue.
\item[\texttt{retrieved\_on}] Date de mise dans le projet. Vide si le fichier n'est pas là.
\item[\texttt{resolution}] Résolution déclarée, ou sa limite.
\item[\texttt{notes}] Ce que le fichier ne permet pas de dire.
\item[\texttt{status}] Les quatre fichiers intégrés au prototype sont \texttt{READY}.
\end{description}

État actuel : la carte couleur, le MNT de la fenêtre de calcul, la grille dérivée et l'image de relief correspondante sont intégrés. Les fonctionnalités optionnelles sans données (géologie, galerie d'images et météo) ont été retirées de l'interface.

\section{Rôle de chaque fichier}

\begin{description}
\item[\texttt{index.html}] Page d'accueil. Carte de contexte, lecture des coordonnées, distance à vol d'oiseau, lien vers le calculateur, mode présentation.
\item[\texttt{style.css}] Mise en page sombre, panneau de carte clair, téléphone et ordinateur.
\item[\texttt{script.js}] Charge les JSON, place les clics, calcule \(d\), refuse de tracer une route sur cette image.
\item[\texttt{route-demo.html}] Page de calcul. Contexte à gauche, terrain et chemins à droite, trois fiches.
\item[\texttt{route-demo.css}] Style de cette page seulement.
\item[\texttt{route-demo.js}] Points, plafond de pente, appel du pathfinder, dessin SVG, texte des compromis.
\item[\texttt{pathfinder.js}] A* à 8 voisins. Pas d'accès au DOM. Trois poids : court, doux, crête.
\item[\texttt{data/route-grid.js}] Grille numérique utilisée par le calcul.
\item[\texttt{assets/maps/orbital.jpg}] Aperçu JPEG du mosaïque couleur USGS Gale \(25\,\mathrm{cm}\), recadré et réduit.
\item[\texttt{assets/maps/route-terrain.jpg}] Image de terrain de la fenêtre de calcul. Seul fond sur lequel les chemins sont dessinés.
\item[\texttt{data/dem/route\_dem\_source.tif}] MNT source de cette fenêtre. Le navigateur ne le lit pas : il lit la grille déjà préparée.
\item[\texttt{data/dem/README.md}] Limites du MNT : pas d'altitude inventée depuis la photo ; pas de réemploi d'un MNT d'une autre boîte.
\item[\texttt{raw/}] Gros fichiers locaux, ignorés par git. Le TIFF couleur d'origine n'est pas chargé dans le navigateur.
\item[\texttt{tools/check\_manifest.py}] Vérifie qu'un fichier présent a une source, un crédit et une date.
\item[\texttt{tools/clip\_raster.py}] Découpe un raster selon une boîte. Il ne sert pas au calcul dans la page.
\item[\texttt{DATA\_SETUP\_GUIDE.md}] Guide d'acquisition. Il ne décrit pas la démo déjà calculée.
\item[\texttt{README.md}] Comment ouvrir les deux pages.
\item[\texttt{data/route-data-notes.md}] Note courte sur A* et sur ce que le modèle ne contient pas.
\end{description}

\section{Ce que le projet ne calcule pas}

\begin{itemize}
\item Pas de pente lue dans la luminosité de la photo couleur.
\item Pas de route rover certifiée, pas de traction, pas de rochers, pas de patinage.
\item Pas de géologie ou de minéralogie locale dans cette version.
\item Pas d'appel à l'ancienne API Mars Rover Photos, et pas de clé dans le navigateur.
\end{itemize}

\section{Ouvrir le projet}

Dans le dossier \texttt{nasa\_space\_2} :
\begin{verbatim}
python -m http.server 8000
\end{verbatim}
Puis \url{http://localhost:8000/} pour le contexte, et
\url{http://localhost:8000/route-demo.html} pour les trois chemins.
Le bouton \emph{Reset sample points} retrouve l'exemple
\(210{,}6\,\mathrm{m}\), \(221{,}2\,\mathrm{m}\) et \(226{,}4\,\mathrm{m}\).

\end{document}
